\chapter{Analyse des besoins et identification des cas d’usage}

\label{chap:3}

Le benchmark du chapitre précédent a montré ce que les solutions WMS matures savent couvrir ; il ne dit pas, à lui seul, quelles capacités doivent être retenues dans l’architecture cible. Cette étape consiste donc à revenir aux besoins opérationnels, aux contraintes des produits et aux limites d’infrastructure. L’objectif est de transformer les constats en cas d’usage traçables, sans supposer que toutes les entreprises disposent des mêmes équipements. La numérotation des domaines Réception, Stockage \& Inventaire et Préparation \& Expédition reprend strictement celle des diagrammes UML déjà produits pendant le stage (UC-R01 à UC-R12, UC-S01 à UC-S14 et UC-E01 à UC-E15).


% Original : page 18, section 3.1

\section{Approche SCAN, PLAN et ACT appliquée aux besoins}

\label{sec:3.1}

La démarche SCAN \(\rightarrow\) PLAN \(\rightarrow\) ACT a servi de fil conducteur à l’identification des besoins. Dans la phase SCAN, le travail a consisté à observer le fonctionnement d’un WMS, à comparer les solutions leaders, puis à relever les points de friction possibles : temps d’attente, erreurs de saisie, écarts de stock, ruptures en zone picking, congestion, manque de visibilité ou contraintes produits mal prises en compte. Le cas industriel étudié dans la littérature sur la réception assistée par vision confirme l’intérêt de partir du processus réel, de mesurer les opérations manuelles et de valider ensuite le gain d’une automatisation ciblée plutôt que d’introduire l’IA sans problème clairement défini \cite{ref11}.


Dans la phase PLAN, chaque problème est traduit en un cas d’usage, puis classé selon trois niveaux. Le niveau Core correspond aux fonctions nécessaires au fonctionnement du WMS. Le niveau Smart regroupe les capacités de prédiction, optimisation ou détection qui améliorent la décision sans exiger nécessairement un nouvel équipement physique. Enfin, le niveau Optional désigne les capacités dépendantes d’une infrastructure particulière : caméra 3D, RFID, AMR/AGV, Pick-to-Light, WCS/WES ou BMS/EMS. La phase ACT, dans le périmètre de ce stage de conception, ne correspond pas à un déploiement industriel ; elle correspond à la formalisation UML, à la vérification des scénarios et aux itérations de validation avec l’encadrement.


\reportfigure{scan_plan_act}{Démarche SCAN, PLAN et ACT appliquée au stage de conception.}

% Original : page 18, section 3.2

\section{Problèmes et contraintes identifiés par domaine}

\label{sec:3.2}

L’analyse des besoins a fait apparaître des problèmes de nature différente. Certains sont directement liés au flux physique, comme la congestion à la réception ou les déplacements inutiles. D’autres concernent la qualité de l’information, par exemple un stock théorique différent du stock réel. D’autres encore sont liés aux contraintes produits, à l’environnement, à l’énergie ou aux ressources humaines. Cette diversité justifie une architecture modulaire : un entrepôt peut avoir besoin d’optimisation du slotting sans robotique, tandis qu’un site pharmaceutique peut prioriser la température, les lots et le FEFO. Les recherches récentes sur les architectures IoT pour entrepôts confirment d’ailleurs l’intérêt d’une approche multi-domaines et soulignent les enjeux de temps réel, d’interopérabilité et de déploiement progressif \cite{ref3}\cite{ref12}.


\begingroup
\small\setlength{\tabcolsep}{4pt}
\setlength{\TableWidth}{\dimexpr\linewidth-24pt\relax}
\begin{longtable}{@{}P{0.180\TableWidth}P{0.290\TableWidth}P{0.250\TableWidth}P{0.280\TableWidth}@{}}
\caption{Problèmes, contraintes et orientations par domaine}\label{tab:problemes}\\
\toprule
\textbf{Domaine} & \textbf{Problèmes / contraintes observés} & \textbf{Conséquence métier} & \textbf{Orientation de réponse} \\
\midrule\endfirsthead
\multicolumn{4}{l}{\small\itshape Tableau \thetable{} (suite)}\\
\toprule
\textbf{Domaine} & \textbf{Problèmes / contraintes observés} & \textbf{Conséquence métier} & \textbf{Orientation de réponse} \\
\midrule\endhead
\midrule\multicolumn{4}{r}{\small\itshape Suite à la page suivante}\\\endfoot
\bottomrule\endlastfoot
Réception & Attente quai, comptage manuel, erreurs d’étiquette, colis non conformes & Retard, mauvaise entrée stock, contrôle tardif & Planification, OCR/vision, caméra 3D, gestion par exception \\[3pt]
Stockage \& inventaire & Mauvais emplacement, écarts stock, rupture picking, congestion & Déplacements, blocage préparation, faible fiabilité & Slotting, cycle count ciblé, prédiction rupture, RFID \\[3pt]
Commandes / expédition & Priorités mal gérées, parcours longs, erreurs picking/packing & Retards et baisse du taux de service & Wave, allocation, optimisation parcours, contrôles \\[3pt]
Yard \& quais & Files camion, portes occupées, manque de visibilité transporteur & Temps d’attente et saturation des quais & Rendez-vous, affectation dynamique, suivi cour \\[3pt]
Produits sensibles & Température, humidité, fragilité, lot, DLC & Risque qualité, perte produit, non-conformité & Profil produit, capteurs, FEFO, alerte \\[3pt]
Ressources / énergie & Charge déséquilibrée, équipement indisponible, surconsommation & Sous-productivité, coût, interruption & Affectation, maintenance, WCS/WES, BMS/EMS \\[3pt]
Pilotage & Manque de visibilité et détection tardive & Décision réactive & KPI, prévision, anomalies, recommandations \\[3pt]
\end{longtable}
\endgroup


Cette cartographie des problèmes ne débouche pas systématiquement sur une solution d’intelligence artificielle. Une règle métier, un processus mieux conçu, un simple scanner ou une automatisation classique suffisent souvent. L’intelligence artificielle n’est retenue que lorsque la décision dépend de plusieurs variables croisées, d’un historique, d’une prédiction ou d’une détection difficile à exprimer sous forme de règle simple  :  c’est ce principe qui guide la classification Core / Smart / Optional développée dans les sections suivantes.


% Original : page 19, section 3.3

\section{Cartographie des cas d’usage  :  Réception (UC-R01 à UC-R12)}

\label{sec:3.3}

La réception constitue le premier point de fiabilisation du stock. Le noyau couvre le rendez-vous, la réception, l’identification, les contrôles, la gestion des anomalies et l’orientation vers un emplacement. Les extensions Smart et Optional permettent ensuite d’anticiper la charge et d’automatiser certaines vérifications. Cette logique reste cohérente avec les résultats du cas Bosch, où l’automatisation visuelle de la réception a réduit les temps de traitement mais a aussi révélé des risques de reconnaissance et d’intégration ERP qui nécessitent un contrôle des exceptions \cite{ref11}.


\begingroup
\small\setlength{\tabcolsep}{4pt}
\setlength{\TableWidth}{\dimexpr\linewidth-24pt\relax}
\begin{longtable}{@{}P{0.130\TableWidth}P{0.360\TableWidth}P{0.180\TableWidth}P{0.330\TableWidth}@{}}
\caption{Cas d’usage de réception}\label{tab:reception}\\
\toprule
\textbf{ID} & \textbf{Use case} & \textbf{Type} & \textbf{Condition / moyen principal} \\
\midrule\endfirsthead
\multicolumn{4}{l}{\small\itshape Tableau \thetable{} (suite)}\\
\toprule
\textbf{ID} & \textbf{Use case} & \textbf{Type} & \textbf{Condition / moyen principal} \\
\midrule\endhead
\midrule\multicolumn{4}{r}{\small\itshape Suite à la page suivante}\\\endfoot
\bottomrule\endlastfoot
UC-R01 & Planifier rendez-vous et quai & Core & Planning transport / TMS-YMS \\[3pt]
UC-R02 & Réceptionner les marchandises & Core & Flux inbound \\[3pt]
UC-R03 & Identifier le produit & Core & Scan / référence / données de référence \\[3pt]
UC-R04 & Contrôler la quantité & Core & Contrôle attendu vs reçu \\[3pt]
UC-R05 & Contrôler la qualité & Core & Règles qualité / opérateur \\[3pt]
UC-R06 & Gérer une anomalie de réception & Core & Écart ou non-conformité \\[3pt]
UC-R07 & Proposer emplacement de stockage & Core & Règles capacité / compatibilité \\[3pt]
UC-R08 & Prévoir la charge de réception & Smart & Historique, ASN, rendez-vous, retards \\[3pt]
UC-R09 & Compter les cartons par caméra 3D & Optional & Caméra 3D disponible \\[3pt]
UC-R10 & Lire l’étiquette par OCR & Optional & Caméra + OCR \\[3pt]
UC-R11 & Détecter les défauts par vision & Optional & Caméra + modèle vision \\[3pt]
UC-R12 & Suivre le conteneur connecté & Optional & Capteurs / IoT / RFID \\[3pt]
\end{longtable}
\endgroup


% Original : page 20, section 3.4

\section{Cartographie des cas d’usage  :  Stockage \& Inventaire (UC-S01 à UC-S14)}

\label{sec:3.4}

Le stockage doit garantir deux objectifs simultanés : placer le produit au bon endroit et maintenir une image fiable du stock. Les fonctions Core couvrent donc le putaway, les stocks, les emplacements, l’inventaire, le réapprovisionnement ainsi que les contraintes de lots et de produits sensibles. Les fonctions Smart cherchent à réduire les déplacements, cibler les contrôles et anticiper les ruptures. Les moyens RFID ou robotisés restent optionnels et ne doivent pas être considérés comme une condition préalable au Smart WMS.


\begingroup
\small\setlength{\tabcolsep}{4pt}
\setlength{\TableWidth}{\dimexpr\linewidth-24pt\relax}
\begin{longtable}{@{}P{0.130\TableWidth}P{0.360\TableWidth}P{0.180\TableWidth}P{0.330\TableWidth}@{}}
\caption{Cas d’usage de stockage et d’inventaire}\label{tab:stockage}\\
\toprule
\textbf{ID} & \textbf{Use case} & \textbf{Type} & \textbf{Condition / moyen principal} \\
\midrule\endfirsthead
\multicolumn{4}{l}{\small\itshape Tableau \thetable{} (suite)}\\
\toprule
\textbf{ID} & \textbf{Use case} & \textbf{Type} & \textbf{Condition / moyen principal} \\
\midrule\endhead
\midrule\multicolumn{4}{r}{\small\itshape Suite à la page suivante}\\\endfoot
\bottomrule\endlastfoot
UC-S01 & Mettre en stock (Putaway) & Core & Tâche de mise en stock \\[3pt]
UC-S02 & Gérer les stocks & Core & Mouvements / statuts / quantités \\[3pt]
UC-S03 & Gérer les emplacements & Core & Zones, capacité, disponibilité \\[3pt]
UC-S04 & Réaliser l’inventaire tournant & Core & Cycle count \\[3pt]
UC-S05 & Réapprovisionner la zone picking & Core & Seuils / demande picking \\[3pt]
UC-S06 & Détecter la congestion d’une zone & Smart & Tâches ouvertes, temps d’attente, occupation \\[3pt]
UC-S07 & Surveiller les produits sensibles & Core conditionnel & Profil produit \\[3pt]
UC-S08 & Gérer lots / DLC / FEFO & Core conditionnel & Produit loti ou périssable \\[3pt]
UC-S09 & Choisir l’emplacement optimal (Slotting) & Smart & Rotation, poids, volume, occupation \\[3pt]
UC-S10 & Cibler les emplacements à recompter & Smart & Historique écarts / risque \\[3pt]
UC-S11 & Prédire une rupture avant blocage & Smart & Stock, consommation, commandes \\[3pt]
UC-S12 & Chercher un emplacement alternatif & Smart & Écart ou article introuvable \\[3pt]
UC-S13 & Suivre les palettes par RFID & Optional & Infrastructure RFID \\[3pt]
UC-S14 & Déplacer le stock par AMR / AGV & Optional & Robot + WCS/WES \\[3pt]
\end{longtable}
\endgroup


% Original : page 20, section 3.5

\section{Cartographie des cas d’usage  :  Commandes \& Exécution / Expédition (UC-E01 à UC-E15)}

\label{sec:3.5}

Le flux outbound traduit la demande client en travail exécutable. La chaîne Core va de la réception de la commande à l’expédition, en passant par la planification, la wave, l’allocation, la création des tâches, le picking, le packing et le chargement. Les extensions Smart portent sur la détection de retard, l’optimisation du parcours et la recommandation d’emballage. Les contrôles par vision, pesage ou Pick-to-Light sont optionnels car ils dépendent de l’équipement présent sur le site.


\begingroup
\small\setlength{\tabcolsep}{4pt}
\setlength{\TableWidth}{\dimexpr\linewidth-24pt\relax}
\begin{longtable}{@{}P{0.130\TableWidth}P{0.360\TableWidth}P{0.180\TableWidth}P{0.330\TableWidth}@{}}
\caption{Cas d’usage de commandes et d’expédition}\label{tab:expedition}\\
\toprule
\textbf{ID} & \textbf{Use case} & \textbf{Type} & \textbf{Condition / moyen principal} \\
\midrule\endfirsthead
\multicolumn{4}{l}{\small\itshape Tableau \thetable{} (suite)}\\
\toprule
\textbf{ID} & \textbf{Use case} & \textbf{Type} & \textbf{Condition / moyen principal} \\
\midrule\endhead
\midrule\multicolumn{4}{r}{\small\itshape Suite à la page suivante}\\\endfoot
\bottomrule\endlastfoot
UC-E01 & Recevoir la commande client & Core & ERP / OMS \\[3pt]
UC-E02 & Planifier la commande & Core & Priorité / cut-off / capacité \\[3pt]
UC-E03 & Créer la wave / batch & Core & Regroupement des ordres \\[3pt]
UC-E04 & Allouer le stock & Core & Disponibilité / règles allocation \\[3pt]
UC-E05 & Créer les tâches de picking & Core & Moteur de tâches \\[3pt]
UC-E06 & Effectuer le picking & Core & Opérateur / équipement \\[3pt]
UC-E07 & Effectuer le packing & Core & Contrôle et conditionnement \\[3pt]
UC-E08 & Charger le camion & Core & Quai / transport \\[3pt]
UC-E09 & Expédier la commande & Core & Confirmation sortie \\[3pt]
UC-E10 & Guider le picking RF / Pick-to-Light & Optional & Terminal RF / Pick-to-Light \\[3pt]
UC-E11 & Optimiser le parcours de picking & Smart & Carte, tâches, contraintes allées \\[3pt]
UC-E12 & Contrôler le picking scan / poids / vision & Optional & Scanner / balance / caméra \\[3pt]
UC-E13 & Recommander l’emballage & Smart & Dimensions, poids, fragilité \\[3pt]
UC-E14 & Contrôler le chargement scan / caméra & Optional & Scanner / caméra quai \\[3pt]
UC-E15 & Détecter commande à risque de retard & Smart & Avancement, charge, cut-off \\[3pt]
\end{longtable}
\endgroup


% Original : page 21, section 3.6

\section{Cartographie des cas d’usage  :  Yard \& Quais}

\label{sec:3.6}

La cour (yard) et les quais constituent un domaine explicite car ils conditionnent directement la fluidité de la réception et de l’expédition. Les identifiants UC-Y01 à UC-Y07 suivent les conventions du projet et doivent être repris dans les futurs diagrammes UML. Le domaine reste indépendant d’un YMS avancé : les fonctions de base peuvent être gérées par le WMS, tandis que le suivi temps réel ou la réaffectation dynamique nécessitent davantage de données ou une intégration TMS/YMS.


\begingroup
\small\setlength{\tabcolsep}{4pt}
\setlength{\TableWidth}{\dimexpr\linewidth-24pt\relax}
\begin{longtable}{@{}P{0.130\TableWidth}P{0.360\TableWidth}P{0.180\TableWidth}P{0.330\TableWidth}@{}}
\caption{Cas d’usage de gestion de la cour et des quais}\label{tab:yard}\\
\toprule
\textbf{ID proposé} & \textbf{Use case} & \textbf{Type} & \textbf{Critère d’activation} \\
\midrule\endfirsthead
\multicolumn{4}{l}{\small\itshape Tableau \thetable{} (suite)}\\
\toprule
\textbf{ID proposé} & \textbf{Use case} & \textbf{Type} & \textbf{Critère d’activation} \\
\midrule\endhead
\midrule\multicolumn{4}{r}{\small\itshape Suite à la page suivante}\\\endfoot
\bottomrule\endlastfoot
UC-Y01 & Planifier les rendez-vous transporteurs & Core & Planning livraison / enlèvement \\[3pt]
UC-Y02 & Enregistrer entrée / sortie véhicule & Core & Contrôle accès / quai \\[3pt]
UC-Y03 & Affecter un quai ou une porte & Core & Disponibilité et type de flux \\[3pt]
UC-Y04 & Suivre le statut des quais / portes & Core & État occupé, libre, bloqué \\[3pt]
UC-Y05 & Suivre véhicule / transporteur dans la cour & Optional & YMS, géolocalisation ou check-points \\[3pt]
UC-Y06 & Détecter congestion cour / quais & Smart & Temps d’attente, files, occupation \\[3pt]
UC-Y07 & Recommander une réaffectation dynamique de quai & Smart & Retard, priorité, capacité disponible \\[3pt]
\end{longtable}
\endgroup


% Original : page 21, section 3.7

\section{Cartographie des cas d’usage  :  Ressources \& Infrastructures}

\label{sec:3.7}

Ce domaine regroupe les moyens nécessaires à l’exécution : personnes, équipements, actifs et systèmes de contrôle. Il joue un rôle de support pour tous les processus métier. Les identifiants UC-RES01 à UC-RES10 suivent les conventions du projet. Ils décrivent les besoins de gestion, tandis que les cas énergétiques de la section~\ref{sec:3.8} détaillent la mesure, l’analyse et l’action. Le suivi énergétique peut utiliser des relevés disponibles ; la collecte automatisée et les interfaces WCS/WES ou BMS/EMS restent conditionnées par l’infrastructure. Une recommandation de veille ne signifie pas un arrêt automatique : l’application d’une consigne relève du système externe et de la politique de validation retenue.


\begingroup
\small\setlength{\tabcolsep}{4pt}
\setlength{\TableWidth}{\dimexpr\linewidth-24pt\relax}
\begin{longtable}{@{}P{0.160\TableWidth}P{0.330\TableWidth}P{0.180\TableWidth}P{0.330\TableWidth}@{}}
\caption{Cas d’usage des ressources et infrastructures}\label{tab:ressources}\\
\toprule
\textbf{ID proposé} & \textbf{Use case} & \textbf{Type} & \textbf{Critère d’activation} \\
\midrule\endfirsthead
\multicolumn{4}{l}{\small\itshape Tableau \thetable{} (suite)}\\
\toprule
\textbf{ID proposé} & \textbf{Use case} & \textbf{Type} & \textbf{Critère d’activation} \\
\midrule\endhead
\midrule\multicolumn{4}{r}{\small\itshape Suite à la page suivante}\\\endfoot
\bottomrule\endlastfoot
UC-RES01 & Planifier la main-d’œuvre & Core & Effectifs, horaires, activité prévue \\[3pt]
UC-RES02 & Gérer compétences et habilitations & Core & Profils opérateurs \\[3pt]
UC-RES03 & Suivre performance et charge des équipes & Core & Tâches et temps d’exécution \\[3pt]
UC-RES04 & Optimiser l’équilibrage de charge & Smart & Charge, distance, compétence, disponibilité \\[3pt]
UC-RES05 & Suivre disponibilité / état des équipements & Core & Chariot, convoyeur, robot, AS/RS \\[3pt]
UC-RES06 & Gérer maintenance et indisponibilités & Core & État, alerte, cycle de vie \\[3pt]
UC-RES07 & Orchestrer les équipements via WCS / WES & Optional & WCS/WES et automatisation présents \\[3pt]
UC-RES08 & Suivre actifs et taux d’utilisation & Core & Données équipement / actif \\[3pt]
UC-RES09 & Collecter données énergie / environnement & Optional & Compteurs et capteurs \\[3pt]
UC-RES10 & Déclencher consignes via BMS / EMS & Optional & BMS/EMS disponible \\[3pt]
\end{longtable}
\endgroup


% Original : page 22, section 3.8

\section{Cartographie des cas d’usage  :  Intelligence \& Aide à la décision transversale}

\label{sec:3.8}

Les capacités transversales ont déjà été esquissées dans les diagrammes UML du stage ; leur numérotation est donc conservée. Elles ne remplacent pas les cas d’usage métier, mais les alimentent. Par exemple, UC-AI02 peut déclencher ou recommander UC-S05, tandis que UC-AI03 soutient UC-E15. Les services de température et d’énergie reposent sur des capteurs et peuvent, lorsque l’infrastructure le permet, transmettre une consigne à un BMS/EMS. Cette séparation permet de distinguer clairement la donnée mesurée, l’analyse intelligente et l’action métier.


\begingroup
\small\setlength{\tabcolsep}{4pt}
\setlength{\TableWidth}{\dimexpr\linewidth-24pt\relax}
\begin{longtable}{@{}P{0.220\TableWidth}P{0.220\TableWidth}P{0.360\TableWidth}P{0.200\TableWidth}@{}}
\caption{Familles de cas d’usage transversaux}\label{tab:transversal}\\
\toprule
\textbf{Famille} & \textbf{IDs conservés} & \textbf{Capacités} & \textbf{Type dominant} \\
\midrule\endfirsthead
\multicolumn{4}{l}{\small\itshape Tableau \thetable{} (suite)}\\
\toprule
\textbf{Famille} & \textbf{IDs conservés} & \textbf{Capacités} & \textbf{Type dominant} \\
\midrule\endhead
\midrule\multicolumn{4}{r}{\small\itshape Suite à la page suivante}\\\endfoot
\bottomrule\endlastfoot
Température / environnement & UC-T01 à UC-T05 & Mesurer T°/humidité, détecter seuil, alerter, adapter température & Optional + Smart \\[3pt]
Énergie & UC-EN01 à UC-EN06 & Mesurer consommation, détecter dérive, prévoir, recommander veille, appliquer économie & Optional + Smart \\[3pt]
IA / décision & UC-AI01 à UC-AI06 & Prévoir demande/charge, rupture, retard, anomalies, recommandation, validation humaine & Smart \\[3pt]
\end{longtable}
\endgroup


Dans cette nomenclature, Smart désigne l’aide à la décision au sens large, sans imposer le recours au machine learning. La détection d’un seuil (UC-T03) et l’alerte (UC-T04) peuvent être réalisées par des règles déterministes. La validation humaine (UC-AI06) est un contrôle du processus décisionnel. Les mesures de température, d’humidité et d’énergie dépendent, quant à elles, des sources disponibles. Les catégories Smart et Optional peuvent donc se cumuler.

Cette organisation reste cohérente avec les travaux récents sur les réseaux IoT pour entrepôts, qui mettent en avant le rôle du monitoring temps réel, de l’analytics et de l’intelligence distribuée, tout en rappelant les contraintes d’interopérabilité, de sécurité et de montée en charge \cite{ref12}. Elle est également compatible avec une logique de performance durable, dans laquelle les données d’énergie et d’automatisation peuvent être suivies comme des indicateurs opérationnels et non comme un sous-sujet isolé \cite{ref13}.


% Original : page 22, section 3.9

\section{Synthèse Core / Smart / Optional et critères d’activation}

\label{sec:3.9}

La classification finale ne vise pas à opposer trois familles rigides ; elle sert surtout à éviter deux erreurs : considérer toute fonction avancée comme de l’IA et supposer que tout Smart WMS doit disposer de robotique ou de vision. Une fonction Smart peut fonctionner uniquement avec les données WMS, alors qu’une fonction Optional dépend d’un dispositif matériel ou d’un système externe. Cette distinction constitue le principe de modularité qui guidera l’architecture fonctionnelle du chapitre suivant.


\reportfigure{classification_portrait}{Classification des capacités et vérification de leurs conditions d’activation.}

\begingroup
\small\setlength{\tabcolsep}{4pt}
\setlength{\TableWidth}{\dimexpr\linewidth-32pt\relax}
\begin{longtable}{@{}P{0.170\TableWidth}P{0.230\TableWidth}P{0.220\TableWidth}P{0.220\TableWidth}P{0.160\TableWidth}@{}}
\caption{Classification et conditions d’activation par domaine}\label{tab:classification}\\
\toprule
\textbf{Domaine} & \textbf{Core} & \textbf{Smart} & \textbf{Optional / infrastructure} & \textbf{Critère d’activation} \\
\midrule\endfirsthead
\multicolumn{5}{l}{\small\itshape Tableau \thetable{} (suite)}\\
\toprule
\textbf{Domaine} & \textbf{Core} & \textbf{Smart} & \textbf{Optional / infrastructure} & \textbf{Critère d’activation} \\
\midrule\endhead
\midrule\multicolumn{5}{r}{\small\itshape Suite à la page suivante}\\\endfoot
\bottomrule\endlastfoot
Réception & UC-R01–UC-R07 & UC-R08 & UC-R09–UC-R12 & historique + caméra/OCR/IoT selon cas \\[0pt]
Stockage \& inventaire & UC-S01–UC-S05, UC-S07– UC-S08 & UC-S06, UC-S09–UC-S12 & UC-S13–UC-S14 & données stock ; RFID ou robotique si présente \\[0pt]
Commandes / expédition & UC-E01–UC-E09 & UC-E11, UC-E13, UC-E15 & UC-E10, UC-E12, UC-E14 & données tâches ; RF/P2L, balance ou caméra si présents \\[0pt]
Yard \& quais & UC-Y01–UC-Y04 & UC-Y06–UC-Y07 & UC-Y05 & TMS/YMS et suivi cour si disponible \\[0pt]
Ressources \& infrastructures & UC-RES01–03, UC-RES05– 06, UC-RES08 & UC-RES04 & UC-RES07, UC-RES09–10 & WCS/WES, capteurs ou BMS/EMS \\[0pt]
Transversal & --- & UC-AI01–UC-AI06 ; UC-T03–UC-T04 ; UC-EN03–UC-EN05 & UC-T01–UC-T02, UC-T05 ; UC-EN01–UC-EN02, UC-EN06 & qualité des données + infrastructure terrain \\[0pt]
\end{longtable}
\endgroup


Les fonctions UC-S07 et UC-S08 sont Core conditionnelles : elles deviennent nécessaires lorsque les caractéristiques des produits imposent une surveillance, une gestion des lots ou une règle FEFO. Les moyens RF, scan ou Pick-to-Light de UC-E10, UC-E12 et UC-E14 sont classés selon leur dépendance matérielle ; cette classification ne remet pas en cause le caractère essentiel du contrôle métier lui-même.

% Original : page 23, section 3.10

\Needspace{17\baselineskip}
\section{Mise en perspective pour l’architecture fonctionnelle cible}

\label{sec:3.10}

L’analyse des besoins transforme une liste de technologies en une structure fonctionnelle cohérente. Les cas d’usage Core définissent le socle commun ; les cas d’usage Smart précisent où les données, l’optimisation ou la prédiction améliorent la décision ; les cas d’usage Optional explicitent les dépendances à l’infrastructure. Le chapitre 4 articule ces capacités dans une architecture générique et modulaire : cœur opérationnel, gestion des tâches, ressources et infrastructures, données, services d’intelligence et terminaux de terrain. Les blocs « Données de référence » et « Stockage des données (OLTP / OLAP-DWH) » constituent des supports transversaux, et non des cas d’usage déclenchés par un acteur. Ils assurent la traçabilité, la cohérence métier et l’exploitation analytique. Il ne s’agit donc plus seulement de choisir des fonctionnalités, mais de positionner chaque capacité, de préciser ses données, ses échanges avec les autres systèmes et ses conditions d’activation.

